\subsection{复形的运算}

在集合 A × A 上，这两个运算

\[
\begin{array}{l} (a, b) + (a ^ {\prime}, b ^ {\prime}) = (a + a ^ {\prime}, b + b ^ {\prime}) \\ (a, b) (a ^ {\prime}, b ^ {\prime}) = (a a ^ {\prime}, a b ^ {\prime} + b a ^ {\prime}) \end{array}
\]

定义一个环结构，记为 A(ε)，其单位元为 1 = (1, 0)；嵌入映射 a ↦ (a, 0) = a1 使得 A 可被等同于 A(ε) 的一个子环；模 A(ε) 是自由的，其基为 \{ 1, ε \} 其中 ε = (0, 1)；我们有 ε² = 0，且 ε 在 A(ε) 中是中心的。

当A是交换的时， \(\mathrm{A}(\varepsilon)\) 是 A 上的双数代数（III，第 15 页）。赋予 \(\mathrm{A}(\varepsilon)\) 环分次结构（II，第164页），其中 \(\mathrm{A}(\varepsilon)_{0} = \mathrm{A}.1\) , \(\mathrm{A}(\varepsilon)_{-1} = \mathrm{A}. \varepsilon\) ，且 \(\mathrm{A}(\varepsilon)_{n} = 0\) （当 \(n \neq 0, -1\) ）。显然，在一个集合 C 上给出 A-复形结构，等同于在 C 上给出一个 \(\mathrm{A}(\varepsilon)\) 分次模结构，其中微分 d 对应于位似 \(\varepsilon_{C}\) ；同样地，复形的态射对应于 \(\mathrm{A}(\varepsilon)\) 分次模的次数为 0 的分次同态。因此， \(\mathrm{A}(\varepsilon)\) 分次模与 A-复形的结构种类是等价的（E, IV, 第 9-10 页）。我们将利用这一事实，将分次模理论中的通常概念移植到复形理论中。

对于分次子 A(ε)-模的概念，对应的是子复形的概念：因此，复形 (C, d) 的子复形是 C 的一个分次子模 C'，使得对每个 \(n \in Z\) , \(d_{n}(C_{n}') \subset C_{n-1}'\) ；若用 \(d'\) 表示由 d 诱导的从 C' 到 C' 的分次 A-同态，则 (C', d') 是一个复形结构，称为由 (C, d) 诱导的结构。除非另有明确说明，每个子复形都配备诱导结构。

我们将其留给读者，按照下面的词典，类似地详细阐述复形的商复形、复形的正合列、复形态射的核、余核和像的概念：

分次商 A(ε)-模

= 商复形，

次数为 0 的分次 A(ε)-同态的核、余核、像 \} 复形态射的核、余核、像，

分次A(ε)-模的正合列以及次数为0的分次同态 \} = 复形的正合列。

例如，第 1 号中的典范正合列给出复形的正合列，称为典范的。

\begin{enumerate}
\def\labelenumi{(\Roman{enumi})}
\tightlist
\item
\end{enumerate}

\[
0 \longrightarrow \mathrm{Z} (\mathrm{C}) \longrightarrow \mathrm{C} \stackrel {\delta} {\longrightarrow} \mathrm{B} (\mathrm{C}) (- 1) \longrightarrow 0,\tag{II}
\]

\[
0 \longrightarrow \mathrm{B} (\mathrm{C}) \longrightarrow \mathrm{Z} (\mathrm{C}) \longrightarrow \mathrm{H} (\mathrm{C}) \longrightarrow 0,\tag{III}
\]

\[
0 \longrightarrow \mathrm{B} (\mathrm{C}) \longrightarrow \mathrm{C} \longrightarrow \mathrm{C} / \mathrm{B} (\mathrm{C}) \longrightarrow 0,\tag{IV}
\]

\[
0 \longrightarrow \mathrm{H} (\dot {\mathrm{C}}) \longrightarrow \mathrm{C} / \mathrm{B} (\mathrm{C}) \stackrel {\delta} {\longrightarrow} \mathrm{B} (\mathrm{C}) (- 1) \stackrel {\cdot} {\longrightarrow} 0,\tag{V}
\]

\[
0 \longrightarrow \mathrm{H} (\mathrm{C}) \longrightarrow \mathrm{C} / \mathrm{B} (\mathrm{C}) \longrightarrow \mathrm{Z} (\mathrm{C}) (- 1) \longrightarrow \mathrm{H} (\mathrm{C}) (- 1) \longrightarrow 0.
\]

类似地，可以定义复形的直和、复形的归纳系以及复形的归纳系的归纳极限等概念。

设 \((\mathbf{C}_{i}, d_{i})\) 为一族复形。称此族的积，并记作 \(\prod_{i \in I} (\mathbf{C}_{i}, d_{i})\)，是按如下方式得到的复形 \((\mathbf{C}, d)\)：

\begin{enumerate}
\def\labelenumi{\alph{enumi})}
\item
  对于每个 \(n \in Z\) , \(C_{n}\) 是直积 A-模 \(\prod_{i \in I} (\mathbf{C}_{i})_{n}\)，即给定复形的齐次分量 \((\mathbf{C}_{i})_{n}\) 的直积
\item
  对于每个 \(n \in Z\) , \(d_{n} : C_{n} \to C_{n-1}\) 是A-同态，其分量为 \((d_{i})_{n}\) .
\end{enumerate}

当 I 有限时， \(\prod_{i\in I}(\mathbf{C}_{i},d_{i})\) 等于 \(\bigoplus_{i\in I}(\mathbf{C}_{i},d_{i})\) 。应当注意，一般而言，积复形 \(\prod_{i\in I}(\mathbf{C}_{i},d_{i})\) 的底层 A-模不是直积 A-模 \(\prod_{i\in I}\mathbf{C}_{i}\) .

考虑一个族（相应地，一个滤过归纳系） \((\mathbf{C}_i)_{i\in I}\) 。设 C 为 \(C_i\) 的直和（相应地，归纳极限），并设 \(\alpha_{i}:C_{i}\to C\) 为典范同态。则 \(\mathrm{H}(\alpha_{i}): \mathrm{H}(\mathrm{C}_{i})\to \mathrm{H}(\mathrm{C})\) 定义了从 \(\bigoplus_{i\in I}\mathrm{H}(\mathrm{C}_{i})\) （相应地， \(\varinjlim_{i\in I}\mathrm{H}(\mathrm{C}_{i})\) ）到 \(\mathrm{H}(\mathrm{C})\) 中的一个次数为 0 的、称为典范的分次同态。同样地，典范投影 \(\prod_{i\in I}\mathrm{C}_{i}\to \mathrm{C}_{i}\) 定义了从 \(\mathrm{H}(\prod \mathrm{C}_i)\) 到 \(\prod (\mathrm{H}(\mathrm{C}_i))\) 的一个次数为 0 的、称为典范的分次同态 .

命题1. --- 对于每个复形族 \((\mathbf{C}_{i})_{i\in I}\) ，典范同态

\[
\bigoplus_ {i \in \mathbf {I}} \mathrm{H} (\mathrm{C} _ {i}) \rightarrow \mathrm{H} \left(\bigoplus_ {i \in \mathbf {I}} \mathrm{C} _ {i}\right), \quad \mathrm{H} \left(\prod_ {i \in \mathbf {I}} \mathrm{C} _ {i}\right)\rightarrow \prod_ {i \in \mathbf {I}} \mathrm{H} (\mathrm{C} _ {i})\tag{2}
\]

都是双射的。

对于复形的每个滤过归纳系 \((\mathbf{C}_{i})_{i\in\mathbf{I}}\) ，典范同态

\[
\varinjlim_ {i \in \mathrm{I}} \mathrm{H} (\mathrm{C} _ {i}) \to \mathrm{H} (\varinjlim_ {i \in \mathrm{I}} \mathrm{C} _ {i})\tag{3}
\]

是双射。

这直接由II，第14页，命题7的推论1，第11页，命题5的推论，以及第91页，命题3得出。

